\documentclass[11pt]{article}
\pdfminorversion=7

\usepackage[margin=1in]{geometry}
\usepackage{booktabs}
\usepackage{graphicx}
\usepackage{tabularx}
\usepackage{amsmath}
\usepackage[numbers,sort&compress]{natbib}
\usepackage[hidelinks]{hyperref}
\setlength{\emergencystretch}{2em}
\newcolumntype{Y}{>{\raggedright\arraybackslash}X}
\newcolumntype{L}[1]{>{\raggedright\arraybackslash}p{#1}}

\title{Temporal and Dataset Shift Expose Shortcut Risk in COVID-19 Respiratory-Audio Screening}

\author{Author names and affiliations to be added before submission}
\date{}

\begin{document}
\maketitle

\begin{abstract}
COVID-19 respiratory-audio benchmarks are often interpreted through internal discrimination, but labels in crowdsourced recordings can be entangled with symptoms, recruitment, protocol, devices, geography, and calendar time. We audited Coswara for source-domain multimodal and temporal analyses, and COUGHVID as an independent cough-only external transfer target. A Coswara cough+speech fusion endpoint reached AUROC 0.897 in a participant-disjoint internal split, but selected endpoints fell to AUROC 0.849 under time-stratified validation and 0.698 under strict early-to-late validation. In model-matched cough-only comparisons using the frozen ComParE+IS10 handcrafted feature strategy, Coswara internal AUROC was 0.849-0.868, whereas COUGHVID transfer AUROC was 0.523-0.543. Metadata-only controls reached AUROC 0.964; symptoms-only and demographic/protocol-only controls reached AUROC 0.932 and 0.914. Inverse-propensity sensitivity reduced the unweighted audio AUROC from 0.879 to 0.780-0.807 depending on weight truncation, supporting a qualified residual audio signal rather than an unconfounded biomarker claim. Source-domain respiratory-audio signal can therefore coexist with non-audio shortcut structure, temporal instability, and weak external portability. COVID-19 audio results should be interpreted by validation target, not as a single leaderboard.
\end{abstract}

\section*{Introduction}

Respiratory audio offered an appealing route for remote digital screening during the COVID-19 pandemic: cough, breath, and speech recordings can be collected without clinical imaging, laboratory infrastructure, or in-person assessment. Public datasets such as Coswara, COUGHVID, and COVID-19 Sounds made respiratory-audio research reproducible at a scale that was difficult to achieve in traditional acoustic studies \citep{bhattacharya2023coswara,orlandic2021coughvid,xia2021covid}. These resources also created a methodological challenge. When data are collected through public participation, labels may reflect testing access, symptom burden, recruitment context, device and environment, geography, and the phase of the pandemic as much as disease-related physiology.

Several respiratory-audio studies have shown that cough, breath, or speech recordings can support apparently strong COVID-19 discrimination within a dataset, using feature-engineered classifiers, deep transfer learning, ensemble methods, and transformer-style spectrogram models \citep{pahar2021global,pahar2022transfer,chowdhury2022mcdm,aytekin2024hst,islam2026dndf}. That literature establishes the plausibility of learnable acoustic signal. It does not, by itself, establish that the learned signal is portable across later collection periods or independent datasets. This distinction matters in medical AI because shortcuts can be stable within a benchmark yet fail when case mix, acquisition, prevalence, or label mechanisms change \citep{geirhos2020shortcut,wynants2020prediction,roberts2021common}.

The resulting methodological question is not whether one more architecture can improve an internal score. The more clinically relevant question is whether a respiratory-audio endpoint retains meaning after its validation target changes. A useful COVID-19 audio paper should therefore separate at least four claims: source-domain discrimination, future-period robustness, resistance to measured non-audio confounding, and external portability. These claims require different evidence and cannot be collapsed into one headline metric.

Reliability-focused COVID-audio studies point in the same direction. Realistic audio-based testing work has shown that study design, participant splitting, and sampling assumptions can materially change performance estimates \citep{han2022sounds}. Matched analyses have questioned whether audio classifiers add evidence beyond symptom checkers once measured confounders are considered \citep{coppock2024audio}. Drift-aware cough-audio work has further shown that performance can change across pandemic time \citep{ganitidis2025jmir}. These findings motivate a validation question that is different from architecture ranking: not whether a particular model achieves the largest internal AUROC, but whether benchmark signal remains reliable after participant separation, temporal stress, confounding checks, and external transfer.

We therefore frame this study as a medical-AI reliability audit of COVID-19 respiratory-audio benchmarks. Coswara is used for source-domain multimodal analyses because it contains cough, breath, speech-style recordings, and metadata. COUGHVID is used only as an external cough-only transfer target; it does not validate multimodal fusion because matching breath and speech modalities are not available for this audit. Following reporting principles for clinical prediction models \citep{collins2024tripodai}, we interpret each result according to the validation target it can support: internal source-domain discrimination, temporal robustness, non-audio shortcut risk, or external cough transfer.

\begin{figure}[t]
\centering
\includegraphics[width=0.92\linewidth]{../common_figures/fig_evaluation_ladder.pdf}
\caption{\textbf{Reliability-audit ladder.} Coswara supports source-domain multimodal evaluation and temporal stress testing, whereas COUGHVID is used only for external cough transfer. The audit separates participant-disjoint internal discrimination from stricter questions about calendar shift, metadata shortcut risk, and dataset-source shift.}
\label{fig:evaluation_ladder}
\end{figure}

\section*{Results}

\subsection*{Validation roles determine the claim level}

The audit used a validation ladder rather than a single headline split (Fig.~\ref{fig:evaluation_ladder}). Each tier was assigned a distinct scientific role (Table~\ref{tab:dataset_roles}). Coswara participant-disjoint validation asks whether the source benchmark contains learnable respiratory-audio signal without participant overlap. Time-stratified and strict early-to-late Coswara validation ask whether that source-domain signal remains stable when collection time is made explicit. Metadata-only controls ask whether labels are predictable from non-audio variables. COUGHVID transfer asks a narrower external question: whether Coswara-trained cough models generalize to an independently collected cough corpus.

\begin{table}[t]
\centering
\caption{\textbf{Dataset and protocol roles.} The sample size is the held-out evaluation size for the primary endpoint in each audit tier. The table is claim-scoped: Coswara supports multimodal source-domain and temporal analyses, whereas COUGHVID supports only cough-only external transfer in this manuscript.}
\label{tab:dataset_roles}
\scriptsize
\setlength{\tabcolsep}{2.5pt}
\begin{tabularx}{\linewidth}{L{0.16\linewidth} L{0.25\linewidth} r L{0.25\linewidth} Y}
\toprule
Dataset & Endpoint in this audit & $n$ & Validation role & Claim not supported \\
\midrule
Coswara & Cough+speech fusion & 314 & Participant-disjoint source-domain test & Future-period or external-dataset performance. \\
Coswara & Cough+breath+speech fusion & 431 & Time-stratified participant test & External validation or causal attribution to audio. \\
Coswara & Breath temporal endpoint & 411 & Strict early-to-late temporal test & COUGHVID transfer or multimodal external validation. \\
Coswara & Metadata-only controls & 2862 & Non-audio shortcut audit & Deployable diagnostic model. \\
COUGHVID & Cough-only external transfer & 8331 & Independent dataset-transfer test & Breath, speech, or multimodal fusion validity. \\
\bottomrule
\end{tabularx}
\end{table}

\subsection*{Coswara contains learnable source-domain respiratory-audio signal}

The participant-disjoint internal split produced strong source-domain discrimination. The selected Coswara cough+speech fusion endpoint reached AUROC 0.897, AUPRC 0.863, balanced accuracy 0.825, F1 0.752, sensitivity 0.833, and specificity 0.816 (Table~\ref{tab:main_results}). This result supports the presence of learnable source-domain signal in the benchmark. It does not establish that the same signal is stable over time or portable to another collection process.

\begin{table}[t]
\centering
\caption{\textbf{Primary validation results.} Rows are the selected held-out endpoints for each protocol, not a fixed-model ablation. BA denotes balanced accuracy. COUGHVID is cough-only external transfer; the Coswara multimodal rows should not be interpreted as externally validated fusion.}
\label{tab:main_results}
\scriptsize
\setlength{\tabcolsep}{2pt}
\begin{tabularx}{\linewidth}{L{0.23\linewidth} L{0.25\linewidth} r r r r r r r}
\toprule
Protocol & Selected endpoint & $n$ & AUROC & AUPRC & BA & F1 & Sens. & Spec. \\
\midrule
Internal participant split & Coswara cough+speech fusion & 314 & 0.897 & 0.863 & 0.825 & 0.752 & 0.833 & 0.816 \\
Time-stratified participant split & Coswara cough+breath+speech fusion & 431 & 0.849 & 0.783 & 0.783 & 0.705 & 0.710 & 0.857 \\
Strict early-to-late temporal split & Coswara breath endpoint & 411 & 0.698 & 0.896 & 0.656 & 0.751 & 0.646 & 0.667 \\
COUGHVID external transfer & Coswara-to-COUGHVID cough model & 8331 & 0.543 & 0.040 & 0.510 & 0.058 & 0.084 & 0.936 \\
\bottomrule
\end{tabularx}
\end{table}

\subsection*{Temporal validation exposes fragility}

The internal signal weakened as validation became more time-aware. Under the time-stratified participant split, the selected endpoint reached AUROC 0.849, AUPRC 0.783, balanced accuracy 0.783, and F1 0.705. Under strict early-to-late temporal validation, AUROC fell to 0.698 and balanced accuracy to 0.656, although AUPRC was 0.896 and F1 was 0.751. The high temporal AUPRC should be interpreted with the validation design and class distribution rather than as evidence against drift; AUROC and balanced accuracy both indicate substantially weaker chronological generalization.

\begin{figure}[t]
\centering
\includegraphics[width=0.92\linewidth]{../common_figures/fig_validation_degradation.pdf}
\caption{\textbf{Performance changes with validation severity.} The plotted bars are selected endpoints from each validation tier, so they summarize claim severity rather than a fixed-model ablation. Source-domain participant-disjoint performance was highest, time-stratified validation was lower, strict early-to-late validation was lower by AUROC and balanced accuracy, and COUGHVID cough-only transfer was near chance. Table~\ref{tab:fixed_cough_ladder} gives the modality- and model-family-matched cough-only comparison.}
\label{fig:validation_degradation}
\end{figure}

This temporal pattern changes the interpretation of the internal result. Participant separation is necessary to prevent direct leakage between training and test participants, but it does not remove calendar-linked sampling effects. In pandemic datasets, calendar time can proxy testing policy, variant waves, recruitment, symptom awareness, health-seeking behavior, and recording conditions. A respiratory-audio endpoint that depends on those correlations can look useful in a source-domain split while losing reliability in later recordings.

\subsection*{Metadata-only controls reveal strong non-audio shortcut risk}

The metadata controls were stronger than the audited audio endpoints in the source-domain setting. Full safe metadata reached AUROC 0.964, AUPRC 0.928, balanced accuracy 0.890, and F1 0.849. Symptoms alone reached AUROC 0.932, and demographic/protocol variables alone reached AUROC 0.914 (Table~\ref{tab:metadata}). These are negative controls for interpretation, not candidate screening tools. They show that label information in Coswara is strongly recoverable from non-audio variables, creating a plausible path for audio models to learn correlated recording, symptom, protocol, or participant-context cues rather than disease-specific acoustics.

Inverse-propensity weighting provided a measured-confounding sensitivity check rather than a causal proof. The unweighted quality-weighted audio fusion row reached AUROC 0.879 and AUPRC 0.832 on 318 held-out samples. Under a conservative cap of 2, AUROC remained 0.807 with effective sample size 238.7, mean absolute residual standardized mean difference (SMD) 0.149, and maximum absolute residual SMD 0.724. Under a stronger cap-10, 99th-percentile truncation setting, AUROC was 0.780, AUPRC 0.537, effective sample size 130.4, mean absolute residual SMD 0.118, and maximum absolute residual SMD 0.382 (Table~\ref{tab:ipw}). The drop from 0.879 to 0.780-0.807 is important: it indicates that audio signal remains visible after measured adjustment, but a substantial fraction of the unweighted result is compatible with observational selection and measured imbalance.

\begin{figure}[t]
\centering
\includegraphics[width=0.88\linewidth]{../common_figures/fig_metadata_confounding.pdf}
\caption{\textbf{Metadata shortcut risk.} Metadata-only controls strongly predicted COVID-19 labels in Coswara. This result indicates that source-domain labels are entangled with symptoms, demographics, protocol variables, and collection context; it does not validate metadata as a clinical diagnostic model.}
\label{fig:metadata_confounding}
\end{figure}

\begin{table}[t]
\centering
\caption{\textbf{Metadata-only confounding audit.} Metadata models were evaluated as controls for benchmark structure. High performance in these rows increases concern that respiratory-audio models trained on the same observational process may exploit correlated non-audio structure.}
\label{tab:metadata}
\scriptsize
\setlength{\tabcolsep}{2pt}
\begin{tabularx}{\linewidth}{L{0.31\linewidth} r r r r r r r Y}
\toprule
Control & $n$ & AUROC & AUPRC & BA & F1 & Sens. & Spec. & Interpretation \\
\midrule
Full safe metadata & 2862 & 0.964 & 0.928 & 0.890 & 0.849 & 0.858 & 0.922 & Labels were highly predictable from non-audio variables. \\
Symptoms only & 2862 & 0.932 & 0.898 & 0.912 & 0.876 & 0.893 & 0.930 & Symptom burden carried strong label information. \\
Demographic/protocol only & 2862 & 0.914 & 0.737 & 0.827 & 0.751 & 0.864 & 0.790 & Non-symptom context also carried label information. \\
Remove recording month in strict temporal metadata audit & -- & +0.247 & -- & -- & -- & -- & -- & Calendar information behaved as an unstable temporal shortcut. \\
\bottomrule
\end{tabularx}
\end{table}

\begin{table}[t]
\centering
\caption{\textbf{Inverse-propensity sensitivity analysis.} IPW rows adjust the quality-weighted audio fusion predictions using measured metadata associated with label assignment and collection context. ESS denotes effective sample size; SMD denotes standardized mean difference.}
\label{tab:ipw}
\scriptsize
\setlength{\tabcolsep}{2pt}
\begin{tabularx}{\linewidth}{L{0.28\linewidth} r r r r r r}
\toprule
Weighting setting & AUROC & AUPRC & BA & ESS & Mean abs. SMD & Max abs. SMD \\
\midrule
Unweighted & 0.879 & 0.832 & 0.804 & 318.0 & 0.224 & 1.384 \\
IPW cap 2, q=0.95 & 0.807 & 0.624 & 0.742 & 238.7 & 0.149 & 0.724 \\
IPW cap 5, q=0.95 & 0.793 & 0.573 & 0.731 & 173.2 & 0.129 & 0.523 \\
IPW cap 10, q=0.99 & 0.780 & 0.537 & 0.721 & 130.4 & 0.118 & 0.382 \\
\bottomrule
\end{tabularx}
\end{table}

\subsection*{Cough-only COUGHVID transfer collapses toward chance}

External transfer was much weaker than Coswara source-domain performance. The best Coswara-to-COUGHVID cough-only row reached AUROC 0.543, AUPRC 0.040, balanced accuracy 0.510, and F1 0.058 on 8331 COUGHVID examples (Table~\ref{tab:coughvid_transfer}). The result is near random by AUROC and balanced accuracy, and the AUPRC is low in the low-prevalence external target. This does not invalidate Coswara or COUGHVID as research datasets. It shows that a Coswara-trained cough representation did not behave as a stable, portable COVID-19 cough marker under direct transfer to a different collection process.

This transfer experiment deliberately uses the frozen ComParE+IS10 handcrafted feature strategy rather than a WavLM or CNN branch. That choice makes the endpoint a standardized stress test of a widely used paralinguistic-feature paradigm under maximum dataset shift, not a complete statement about every possible deep representation. The practical WavLM and CNN-BiGRU branches were source-domain, single-stream model-development checks; they were not completed as COUGHVID transfer experiments and should not be cited as external validation evidence.

The selected-endpoint validation ladder is useful for claim control, but it is not a fixed-model ablation because the selected modality and fusion rule change across tiers. The fair external-transfer comparison is cough-only on both sides, not multimodal Coswara fusion against cough-only COUGHVID. Table~\ref{tab:cough_internal_external} reports this internal-to-external comparison. Under the same ComParE+IS10 top-800 feature strategy, Coswara cough-only internal AUROC ranged from 0.849 to 0.868, whereas the corresponding COUGHVID AUROC ranged from 0.523 to 0.543. Table~\ref{tab:fixed_cough_ladder} extends this check across the full validation ladder while holding modality fixed to cough and model family fixed. The external result is therefore not explained merely by the absence of breath and speech in COUGHVID; even the matched cough branch transported poorly.

\begin{table}[t]
\centering
\caption{\textbf{COUGHVID external-transfer results.} Coswara-trained cough models were evaluated on COUGHVID without pooling COUGHVID into source-domain model selection. These rows test cough-only dataset transfer and should not be read as multimodal external validation.}
\label{tab:coughvid_transfer}
\scriptsize
\setlength{\tabcolsep}{2pt}
\begin{tabularx}{\linewidth}{L{0.30\linewidth} r r r r r r r}
\toprule
Model & $n$ & AUROC & AUPRC & BA & F1 & Sens. & Spec. \\
\midrule
LightGBM with class balancing & 8331 & 0.543 & 0.040 & 0.510 & 0.058 & 0.084 & 0.936 \\
XGBoost with class balancing & 8331 & 0.532 & 0.039 & 0.505 & 0.063 & 0.267 & 0.743 \\
CatBoost with class balancing & 8331 & 0.531 & 0.040 & 0.506 & 0.059 & 0.144 & 0.869 \\
RBF support-vector classifier & 8331 & 0.523 & 0.037 & 0.500 & 0.058 & 0.193 & 0.807 \\
\bottomrule
\end{tabularx}
\end{table}

\begin{table}[t]
\centering
\caption{\textbf{Model-matched cough-only internal-to-external comparison.} Internal rows use the Coswara participant-disjoint cough endpoint with ComParE+IS10 top-800 features. External rows apply the corresponding Coswara-trained cough model to COUGHVID. This table is the modality-matched comparison; it should be read separately from Coswara multimodal fusion.}
\label{tab:cough_internal_external}
\begin{tabular}{lrrrrrr}
\toprule
Model & \multicolumn{3}{c}{Coswara cough-only} & \multicolumn{3}{c}{COUGHVID cough-only} \\
\cmidrule(lr){2-4}\cmidrule(lr){5-7}
 & AUROC & AUPRC & BA & AUROC & AUPRC & BA \\
\midrule
LightGBM-SMOTE & 0.853 & 0.797 & 0.782 & 0.543 & 0.040 & 0.510 \\
SVC-RBF & 0.868 & 0.812 & 0.795 & 0.523 & 0.037 & 0.500 \\
CatBoost-SMOTE & 0.849 & 0.788 & 0.776 & 0.531 & 0.040 & 0.506 \\
XGBoost-SMOTE & 0.850 & 0.780 & 0.776 & 0.532 & 0.039 & 0.505 \\
\bottomrule
\end{tabular}
\end{table}

\begin{table}[t]
\centering
\caption{\textbf{Fixed cough-only model-family ladder.} Each row keeps the modality fixed to cough and follows the same model family across validation tiers. Models are retrained within the corresponding source-domain protocol; values are AUROC.}
\label{tab:fixed_cough_ladder}
\begin{tabular}{lrrrr}
\toprule
Model & Existing split & Time-stratified & Strict temporal & COUGHVID \\
\midrule
LightGBM-SMOTE & 0.853 & 0.830 & 0.604 & 0.543 \\
SVC-RBF & 0.868 & 0.802 & 0.561 & 0.523 \\
CatBoost-SMOTE & 0.849 & 0.826 & 0.614 & 0.531 \\
XGBoost-SMOTE & 0.850 & 0.828 & 0.599 & 0.532 \\
\bottomrule
\end{tabular}
\end{table}

\subsection*{Evidence triangulation separates feasibility from reliability}

The audit is strongest when the results are read together rather than as isolated scores (Table~\ref{tab:evidence_triangulation}). The internal Coswara result shows feasibility: respiratory recordings and acoustic features contain label-associated signal inside the source collection. The temporal results show instability: that signal weakens when calendar structure is made part of validation. The metadata controls show shortcut risk: non-audio variables alone solve much of the source-domain label prediction problem. The COUGHVID result shows weak portability: a Coswara-trained cough rule does not transfer robustly to an independent cough corpus. These findings do not contradict each other. They describe a benchmark in which disease-related acoustic information may be present, but is mixed with collection, symptom, and dataset-source structure.

\begin{table}[t]
\centering
\caption{\textbf{Evidence triangulation and claim boundary.} Each audit component answers a different question. The manuscript conclusion is based on the pattern across components rather than on a single favorable or unfavorable metric.}
\label{tab:evidence_triangulation}
\scriptsize
\setlength{\tabcolsep}{2.5pt}
\begin{tabularx}{\linewidth}{L{0.25\linewidth} L{0.24\linewidth} L{0.25\linewidth} Y}
\toprule
Audit component & Key result & Supports & Does not support \\
\midrule
Participant-disjoint Coswara fusion & AUROC 0.897 & Source-domain respiratory-audio feasibility. & Future-period or external-dataset reliability. \\
Time-aware Coswara validation & AUROC 0.849 and strict temporal AUROC 0.698 & Calendar-sensitive degradation is material. & Prospective clinical deployment. \\
Metadata-only controls & Full metadata AUROC 0.964; symptoms-only AUROC 0.932; demographic/protocol AUROC 0.914 & Strong non-audio label structure and shortcut risk. & Proof of the exact shortcut used by each audio model. \\
Confounding sensitivity analyses & IPW reduced audio AUROC from 0.879 unweighted to 0.807 at cap 2 and 0.780 at cap 10/q=0.99, with ESS 238.7 and 130.4 respectively. & Measured-confounder adjustment is necessary and changes interpretation. & Complete removal of unmeasured confounding or a causal audio effect. \\
COUGHVID cough transfer & AUROC 0.543; AUPRC 0.040 on 8331 samples & Weak cross-dataset cough portability under direct transfer. & Multimodal external validation, because COUGHVID is cough-only here. \\
\bottomrule
\end{tabularx}
\end{table}

A source-domain model-development interpretation would foreground the internal AUROC 0.897. A reliability-audit interpretation asks whether the signal survives plausible deployment shifts: later collection periods, changed prevalence, independent acquisition, and confounding by symptoms or protocol. Under that interpretation, internal discrimination is necessary but not sufficient evidence for clinical screening utility.

\section*{Discussion}

This reliability audit shows that COVID-19 respiratory-audio performance is inseparable from validation design. The selected Coswara internal endpoint achieved AUROC 0.897, but selected endpoints declined to AUROC 0.849 under time-stratified validation and AUROC 0.698 under strict early-to-late temporal validation. In modality-matched cough-only comparisons, Coswara internal AUROC of 0.849-0.868 contrasted with COUGHVID transfer AUROC of 0.523-0.543. Metadata-only controls reached AUROC values above 0.9. These results are not minor implementation details. They indicate that source-domain respiratory-audio discrimination can coexist with shortcut learning, temporal drift, and dataset shift.

The first implication is that participant separation, while essential, is not enough for claims about future-period performance. Pandemic data were collected while testing access, public behavior, symptom reporting, circulating variants, and recruitment routes were changing. Calendar variables can therefore stabilize an internal split while destabilizing later-period evaluation. The metadata ablation is consistent with this interpretation: removing recording month changed strict temporal metadata AUROC by +0.247, suggesting that month was not a reliable disease feature but an unstable proxy for collection context.

The second implication is that non-audio controls should be treated as central evidence in respiratory-audio studies. Symptoms-only AUROC 0.932 and demographic/protocol-only AUROC 0.914 show that labels were strongly recoverable without audio. This does not prove which latent variables the audio models used, but it limits causal interpretation. Recording quality, participant effort, symptom severity, device type, and protocol compliance can all affect respiratory audio and may correlate with infection labels in observational datasets. A model that performs well internally may therefore be detecting a mixture of disease physiology and benchmark-specific context.

The external transfer result defines the strongest boundary of the present evidence. COUGHVID was deliberately used only for cough transfer because it does not provide the matched breath and speech modalities needed to test Coswara multimodal fusion. The near-chance transfer result should therefore be read narrowly and seriously: it is not a failure of multimodal fusion or a completed test of all deep representations, but it is evidence that a standardized Coswara-trained handcrafted cough pipeline did not generalize robustly to an independent cough dataset. This finding is consistent with broader medical-AI concerns about distribution shift and shortcut learning \citep{geirhos2020shortcut,roberts2021common}.

The related literature is best reconciled by separating feasibility from reliability. Prior model-development studies demonstrate that respiratory audio can contain discriminative information under particular dataset and split assumptions \citep{pahar2021global,pahar2022transfer,chowdhury2022mcdm,aytekin2024hst,islam2026dndf}. Reliability-oriented studies show that this information can be sensitive to sampling, symptoms, and time \citep{han2022sounds,coppock2024audio,ganitidis2025jmir}. Our contribution is to place these concerns into one audit ladder for Coswara and COUGHVID: participant-disjoint source-domain testing, temporal stress testing, metadata-only controls, and external cough transfer are evaluated together and interpreted according to their claim level.

The study should therefore not be read as an architecture leaderboard. Transformer spectrogram models, deep decision forests, transfer-learning bottleneck features, and symptom-assisted systems answer important model-development questions. This manuscript asks a different question: which parts of the evidence survive when the validation target changes. That question is directly relevant to digital-medicine translation because a deployable screening system must withstand distribution shift, prevalence shift, threshold transport, and measured confounding rather than only a favorable internal split.

\subsection*{Validity boundaries and limitations}

The limitations of this audit define the boundary of its claims. First, COUGHVID is a cough-only external target in this analysis. It is a valid test of cough transfer from Coswara to an independent cough corpus, but it cannot validate Coswara breath, speech, or multimodal fusion. The correct claim is therefore asymmetric: multimodal Coswara fusion is internally and temporally audited, while external validation is cough-only.

Second, the primary validation ladder reports point estimates rather than structural confidence intervals for every selected endpoint. Paper-comparable cross-validation rows include fold-level variability, but the main participant, temporal, and external-transfer endpoints should ideally be accompanied by DeLong, bootstrap, or paired-resampling confidence intervals. The observed drops are large enough to motivate caution, but small differences between individual models or adjacent protocols should not be interpreted as statistically resolved without those intervals.

Third, both datasets are retrospective and crowdsourced. Their labels, symptoms, devices, environments, recording prompts, recruitment pathways, and population composition are not standardized to the level expected in a prospective clinical diagnostic study. Temporal validation across observed calendar periods is a stricter stress test than a random split, but it remains observational; it cannot replace prospective evaluation in a prespecified screening population.

Fourth, metadata and inverse-propensity analyses identify measured shortcut risk but do not establish a complete causal graph. High symptoms-only and demographic/protocol-only AUROC values show that non-audio label information is present, and residual standardized mean differences after weighting show that adjustment is imperfect. The IPW table should be interpreted as sensitivity evidence: cap-10 weighting improves measured balance but reduces ESS to 130.4, so the retained AUROC 0.780 supports qualified residual signal rather than a definitive causal audio effect. These analyses support conservative interpretation; they do not prove that every audio model used the same shortcut or that all acoustic signal is non-causal.

Fifth, the model bank is broad but not exhaustive. We evaluated engineered acoustic descriptors, boosted and kernel classifiers, fusion, CNN-style branches, self-supervised representations, and feature-search variants, but we did not claim to have exhausted every possible transformer, adaptation, segmentation, or prospective calibration strategy. Additional architectures may improve source-domain performance. The unresolved question for clinical translation is whether they also improve temporally separated and external validation without relying on non-audio shortcuts.

Finally, thresholded metrics and calibration should not be interpreted as clinical operating points. Thresholds were selected within validation protocols, and the external COUGHVID target had low positive prevalence. AUPRC, sensitivity, specificity, and probability calibration are therefore population-dependent. Clinical use would require a prospectively specified target population, recording protocol, threshold policy, calibration plan, and subgroup-safety analysis.

These boundaries strengthen rather than weaken the conclusion: COVID-19 respiratory-audio benchmarks remain valuable for methodological research, but high internal performance should be reported as source-domain discrimination unless supported by temporally separated, externally validated, and confounding-aware evidence. For digital respiratory screening, the burden of proof should rest on reliability under shift rather than internal discrimination alone.

\section*{Methods}

\subsection*{Datasets}

Coswara was used as the source dataset because it contains cough, breathing, speech-style recordings, and participant metadata for SARS-CoV-2 status research \citep{bhattacharya2023coswara}. COUGHVID was used as an external cough-only target dataset \citep{orlandic2021coughvid}. COUGHVID was not pooled into Coswara model development or model selection. Because the external target did not provide matching breath and speech modalities for this audit, COUGHVID transfer was limited to cough models.

\subsection*{Cohort construction and leakage control}

Supervised analyses used known positive and known negative labels. Unknown-label recordings were excluded from supervised performance estimation. Participant identifiers were the leakage-control unit in source-domain splits, so recordings from the same participant did not appear in both training and held-out evaluation within a protocol. Audio recordings were screened for audibility and basic recording integrity before supervised evaluation.

\subsection*{Audio representations and candidate models}

Audio representations included compact spectral and cepstral summaries, OpenSMILE acoustic descriptors, time-frequency representations for convolutional models, and self-supervised speech representations \citep{eyben2010opensmile,schuller2010paralinguistic,schuller2016compare,chen2022wavlm}. Candidate classifiers included regularized linear models, random forests, support-vector machines, LightGBM, XGBoost, and CatBoost \citep{breiman2001random,cortes1995support,ke2017lightgbm,chen2016xgboost,prokhorenkova2018catboost}. Class imbalance handling, where used, was applied within the training data, including synthetic minority oversampling for selected classical learners \citep{chawla2002smote}. Feature selection, including validation-ranked feature subsets and particle-swarm search branches, was performed within the corresponding training/validation workflow \citep{kennedy1995pso}. The final ComParE+IS10 top-800 feature strategy was selected from the Coswara source-domain training workflow and then frozen for downstream validation; COUGHVID labels were not used for feature selection. ComParE+IS10 was used for the transfer endpoint because it provides a standardized, reproducible handcrafted representation comparable to much of the pre-transformer respiratory-audio literature and does not require target-domain fine-tuning. For the strict chronological analysis, this should be read as a frozen source-domain representation rather than a fully prospective feature-selection protocol re-estimated only inside the early-period cohort. Held-out test data were not used for model choice or threshold selection.

\subsection*{Multimodal fusion}

For Coswara source-domain analyses, modality-specific predictions were combined using validation-selected fusion rules, including uniform averaging, validation-weighted averaging, and stacked logistic fusion. Fusion was evaluated only where the required modalities were available. COUGHVID transfer was cough-only and therefore did not evaluate multimodal fusion.

\subsection*{Validation protocols}

The internal participant split tested source-domain discrimination under participant separation. The time-stratified participant split retained participant separation while making collection-time structure explicit. Strict early-to-late temporal validation trained and selected models on earlier recordings and evaluated them on later recordings. The COUGHVID transfer test evaluated Coswara-trained cough models on an independent cough dataset. Each protocol was interpreted according to its validation target rather than as an interchangeable benchmark score.

\subsection*{Metadata confounding audit}

Metadata-only models were trained as controls for non-audio shortcut risk. Feature groups included full safe metadata, symptoms-only metadata, and demographic/protocol-only metadata. These controls were interpreted as evidence about benchmark structure, not as deployable predictors. A temporal ablation removed recording-month information to evaluate whether calendar proxies supported or destabilized strict temporal metadata performance.

\subsection*{Measured-confounding sensitivity}

As a sensitivity analysis, audio predictions were also examined under inverse-propensity weighting using measured metadata associated with label assignment and collection context. The purpose was not to estimate a causal effect of COVID-19 on audio, but to test whether the audio result remained visible after down-weighting measured imbalance. Effective sample size and residual standardized mean differences were inspected because aggressive weighting can create unstable estimates and because measured adjustment cannot remove unmeasured confounding.

\subsection*{Metrics}

We report AUROC, AUPRC, balanced accuracy, F1, sensitivity, and specificity for the primary audit endpoints. AUROC summarizes ranking discrimination across thresholds. AUPRC is prevalence-sensitive and was interpreted with the validation split and class distribution. Balanced accuracy was computed as
\[
\mathrm{BA}=\frac{1}{2}\left(\frac{\mathrm{TP}}{\mathrm{TP}+\mathrm{FN}}+\frac{\mathrm{TN}}{\mathrm{TN}+\mathrm{FP}}\right),
\]
and F1 was computed as
\[
\mathrm{F1}=\frac{2\mathrm{TP}}{2\mathrm{TP}+\mathrm{FP}+\mathrm{FN}}.
\]
Thresholded metrics used thresholds selected from validation data within the corresponding protocol.

\section*{Data availability}

The analysis uses public Coswara and COUGHVID data, available from the dataset maintainers subject to their access terms. Aggregated result summaries and figure source data for the manuscript can be supplied as supplementary material. Individual-level data redistribution should follow the source dataset licenses and participant-consent constraints.

\section*{Code availability}

The audit code will be made available in a public code archive at submission or publication, with documentation sufficient to reproduce the reported aggregate results from the public datasets where licenses permit. During peer review, code can be provided to editors and reviewers under the access conditions required by the journal.

\section*{Acknowledgements}

Acknowledgements and funding information should be completed by the submitting authors.

\section*{Author contributions}

Author contribution statements should be completed with author initials before submission, using standard contribution categories such as conceptualization, methodology, data curation, formal analysis, validation, visualization, writing, supervision, and review.

\section*{Competing interests}

The authors declare no competing interests. This statement should be verified by all authors before submission.

\bibliographystyle{unsrtnat}
\bibliography{references}

\end{document}
